\section{Related Work}
\label{sec:related}

\paragraph{Classical OMR.}
Classical OMR relies on human-designed stages to perform recognition:
binarization, staff removal and detection, symbol segmentation, and musical
reconstruction~\cite{rebelo2012state,calvo2020understanding}. Staff-line
locations are found with projection profiles, run-length encoding, or Hough
transforms~\cite{duda1972hough} and then removed morphologically, after which
symbols can be detected by connected components~\cite{suzuki1985contours}. My
geometry stages are in this tradition; I have built the pipeline with an
engineered phone-capture correction system, with both a front-end rectification
and a post-pipeline stave confirmation by a spacing-consistency check.

\paragraph{Deep and end-to-end OMR.}
Neural sequence models are recent additions to the OMR literature. Calvo-Zaragoza
and Rizo first developed a full end-to-end monophonic OMR CRNN for the PrIMuS
corpus~\cite{calvo2018primus}; van der Wel and Ullrich propose a convolutional
seq-to-seq model for music transcription~\cite{vanderwel2017seq2seq}; Bar\'o et
al.\ use these ideas for handwritten music notation~\cite{baro2019handwritten};
and glyph-based detectors frame glyph locations as boxes~\cite{pacha2018baseline}.
PrIMuS provides the source for the clean symbols in my training data; I do not
use a full sequence model, instead using a CNN~\cite{lecun1998gradient} for the
discrete pitch-identification task to serve as a controlled probe for how
training from clean symbols transfers to photos.

\paragraph{Domain shift.}
A major topic in vision is overcoming the mismatch between training and
deployment domains through unsupervised domain adaptation~\cite{ganin2015dann},
and in OMR this gap has been addressed by creating synthetically distorted
``camera'' versions of the PrIMuS corpus. I introduce an explicit per-class
quantitative measure of this clean-to-real gap in photos, and demonstrate how my
geometry-based pitch-detection approach largely sidesteps it.
